\documentclass[11pt,a4paper]{article}

\usepackage[margin=1in]{geometry}
\usepackage{microtype}
\usepackage{hyperref}
\usepackage{booktabs}
\usepackage{enumitem}
\usepackage{xcolor}
\usepackage{titlesec}
\usepackage{graphicx}
\usepackage{tabularx}
\usepackage[T1]{fontenc}
\usepackage{lmodern}
\usepackage{natbib}
\usepackage{multirow}
\usepackage{array}

\definecolor{accent}{HTML}{2563EB}
\definecolor{codebg}{HTML}{F5F5F5}
\definecolor{codetext}{HTML}{333333}
\definecolor{gray60}{HTML}{666666}

\hypersetup{
  colorlinks=true,
  linkcolor=accent,
  urlcolor=accent,
  citecolor=accent,
}

\titleformat{\section}{\large\bfseries\color{accent}}{}{0em}{}
\titleformat{\subsection}{\normalsize\bfseries}{}{0em}{}
\titleformat{\subsubsection}{\normalsize\itshape}{}{0em}{}

\newcolumntype{L}[1]{>{\raggedright\arraybackslash}p{#1}}
\newcolumntype{C}[1]{>{\centering\arraybackslash}p{#1}}

\title{\textbf{Automated Research Systems}\\[0.3em]
  \large A Survey of Architectures, Approaches, and Trade-Offs\\
  in LLM-Driven Research Automation}
\author{M.\ Berends \and Claude (Anthropic)}
\date{March 2026}

\begin{document}
\maketitle

\begin{abstract}
Large language models have enabled a new class of systems that automate
aspects of the research process---from literature discovery and synthesis
to hypothesis generation, experimental optimisation, and full paper
writing. This survey examines the major approaches to automated research
(``autoresearch''), organising them into four categories: \emph{deep
research} (multi-step web search and report generation), \emph{literature
agents} (scientific question answering over academic corpora),
\emph{full-pipeline science} (end-to-end idea-to-paper systems), and
\emph{empirical optimisation} (autonomous code-level
edit--benchmark--iterate loops). For each category we analyse the
architectures and design decisions of leading systems, identify shared
patterns and key differentiators, and discuss the trade-offs that shape
system behaviour. We find that decomposition strategy, quality-control
mechanisms, and context management are the principal axes along which
systems diverge, and that open-source implementations have reached or
exceeded commercial systems on several benchmarks.
\end{abstract}

% =================================================================
\section{Introduction}
\label{sec:intro}

The application of LLMs to research automation has expanded rapidly
since 2023. Early systems used retrieval-augmented generation (RAG) for
single-turn question answering; current systems employ multi-agent
architectures, iterative search loops, and tool use to conduct extended
autonomous research sessions lasting minutes to hours.

Despite this proliferation, the landscape lacks a unifying framework.
Systems marketed under similar labels---``deep research'',
``AI scientist'', ``research agent''---often automate fundamentally
different aspects of the research process. A literature synthesis agent
that finds and cites papers has little architectural overlap with an
evolutionary code optimiser that discovers new algorithms through
mutation and fitness evaluation.

This survey proposes a four-category taxonomy based on \emph{what is
being automated} (Section~\ref{sec:taxonomy}), then provides detailed
analysis of exemplar systems in each category
(Sections~\ref{sec:deep}--\ref{sec:empirical}). Section~\ref{sec:cross}
identifies cross-cutting patterns, and Section~\ref{sec:conclusion}
discusses implications for future system design.

% =================================================================
\section{Taxonomy}
\label{sec:taxonomy}

We organise automated research systems along two dimensions: the
\emph{input} (what triggers the system) and the \emph{output} (what it
produces). This yields four natural categories:

\begin{table}[h]
\centering
\small
\begin{tabular}{@{}L{2.8cm}L{3.5cm}L{3.5cm}L{4cm}@{}}
\toprule
\textbf{Category} & \textbf{Input} & \textbf{Output} & \textbf{Exemplars} \\
\midrule
Deep Research & Natural-language question & Multi-page cited report &
  OpenAI, Gemini, Perplexity, STORM, GPT-Researcher \\[0.5em]
Literature Agents & Research question & Cited answer with evidence &
  PaperQA2, OpenScholar, Elicit \\[0.5em]
Full-Pipeline Science & Domain + seed ideas & Paper, code, results &
  AI Scientist, ShinkaEvolve, Coscientist, SciAgents \\[0.5em]
Empirical Optimisation & Objective + codebase & Optimised code + log &
  pi-autoresearch, ShinkaEvolve, AlphaEvolve \\
\bottomrule
\end{tabular}
\caption{Taxonomy of automated research systems. ShinkaEvolve appears
in two categories as it combines evolutionary algorithm discovery
(full-pipeline) with code-level optimisation (empirical).}
\label{tab:taxonomy}
\end{table}

\noindent
These categories are not mutually exclusive. The most ambitious systems
(e.g.\ FutureHouse Robin) aim to span multiple categories by combining
literature synthesis with hypothesis generation and experimental
validation.

% =================================================================
\section{Deep Research Systems}
\label{sec:deep}

Deep research systems accept a natural-language question and produce a
comprehensive, cited report after conducting autonomous multi-step web
search. All share a common architectural pattern---the iterative
\emph{ReAct loop} (plan $\to$ search $\to$ read $\to$ reason $\to$
repeat $\to$ synthesise)---but diverge significantly in depth, speed,
and search strategy.

\subsection{Commercial Systems}

\subsubsection{OpenAI Deep Research.}
Launched February 2025 as a ChatGPT feature, OpenAI Deep Research is
powered by a fine-tuned o3-class reasoning model trained via
reinforcement learning on simulated browsing tasks. It executes
30--60 web searches per query across 150--200 reasoning iterations,
visiting dozens of pages over 5--30 minutes. A secondary o3-mini model
summarises the chain-of-thought for user display. At launch it achieved
51.5\% on BrowseComp~\citep{browsecomp2025}---versus 1.9\% for GPT-4o
with browsing---and 26.6\% on Humanity's Last Exam. Its primary
limitation is latency: the system optimises for depth at the expense of
speed, with queries routinely taking 20+ minutes.

\subsubsection{Google Gemini Deep Research.}
Launched December 2024 on Gemini 2.0 Flash Thinking (later Gemini 3
Pro), Google's system introduces an \emph{asynchronous task manager}
that maintains shared state between a planner model and parallel
task-execution models. This enables graceful recovery from individual
sub-task failures without restarting the entire research session.
The user can review and edit the research plan before execution. Native
Google Search integration provides arguably the highest-quality search
backend of any system. On DRACO~\citep{draco2026}, Gemini Deep Research
scores 58.97\%, and on BrowseComp, Gemini-3.1-Pro reaches 85.9.

\subsubsection{Perplexity Deep Research.}
Perplexity optimises for the opposite end of the speed--depth
trade-off. Using a modular planner--retriever--synthesiser architecture
with 3--5 sequential refinement cycles, it completes most queries in
under 3 minutes---approximately 9$\times$ faster than OpenAI. It
references $\sim$50 sources per query on average (vs.\ $\sim$20 for
competitors) via on-demand web crawling rather than static indexes. As
of February 2026, it integrates Claude Opus as the underlying reasoning
model and achieves the highest DRACO score of any system at
70.5\%~\citep{draco2026}. Perplexity also published the DRACO benchmark
itself, designed to evaluate cross-domain deep research quality.

\subsubsection{Grok DeepSearch / DeeperSearch.}
xAI's Grok 4.20 Beta introduced a native 4-agent multi-agent
architecture: a coordinator (Grok), a research specialist (Harper), a
logic/code specialist (Benjamin), and a creative/divergent specialist
(Lucas). All four share the same $\sim$3T parameter MoE backbone with
lightweight LoRA persona adapters. Multi-Agent Reinforcement Learning
(MARL) during post-training rewards convergence on truthful answers.
Uniquely, Grok integrates real-time X/Twitter data, giving it access
to breaking news and social sentiment unavailable to other systems.

\subsubsection{Comparison.}
Table~\ref{tab:deep_commercial} summarises the key parameters.

\begin{table}[h]
\centering
\small
\begin{tabular}{@{}L{2.8cm}C{2cm}C{1.8cm}C{2.2cm}C{2.5cm}@{}}
\toprule
\textbf{System} & \textbf{Searches} & \textbf{Time} & \textbf{BrowseComp} & \textbf{DRACO} \\
\midrule
OpenAI Deep Research & 30--60 & 5--30\,min & 51.5 & 52.1 \\
Gemini Deep Research & 100+ pages & 3--10\,min & 85.9 & 59.0 \\
Perplexity & $\sim$50 sources & 1--3\,min & --- & 70.5 \\
Grok DeepSearch & 10 iterations & 2--3\,min & --- & --- \\
Claude Research & multi-pass & 5--45\,min & 84.0 & --- \\
\bottomrule
\end{tabular}
\caption{Commercial deep research systems. BrowseComp scores use latest
available model variants. Dashes indicate unreported scores.}
\label{tab:deep_commercial}
\end{table}

\subsection{Open-Source Systems}

\subsubsection{GPT-Researcher.}
GPT-Researcher~\citep{gptresearcher2024} uses a plan-and-execute
pattern with a planner agent, researcher agent, and writer agent. The
planner generates sub-queries from the research question; these are
processed \emph{concurrently} via \texttt{asyncio.gather}. Retrieval
supports 12+ backends (Tavily, Google, Bing, DuckDuckGo, Serper, Arxiv,
PubMed, Semantic Scholar, etc.). Scraped content is compressed via
cosine-similarity filtering (threshold 0.35) over 1000-character chunks.

The system's \emph{deep research mode} implements recursive
breadth-first exploration: at each level, it generates \texttt{breadth}
queries (default 4), spawns nested \texttt{GPTResearcher} instances for
each, extracts learnings and follow-up questions via LLM, and recurses
to \texttt{depth} levels (default 2). Context is capped at 25,000 words.
A separate multi-agent mode built on LangGraph adds role-specialised
agents (editor, researcher, writer, reviewer, publisher) with optional
human-in-the-loop review.

\subsubsection{STORM.}
STORM~\citep{storm2024} (Stanford) takes a fundamentally different
approach to question decomposition. Rather than generating sub-queries
directly, it \emph{simulates multi-expert dialogue}:

\begin{enumerate}[nosep]
  \item \textbf{Persona generation.} The system finds related Wikipedia
    topics, scrapes their tables of contents, and asks the LLM to create
    diverse ``editor'' personas (e.g.\ ``Historian focusing on early
    development'', ``Policy analyst examining regulation'').
  \item \textbf{Conversation simulation.} For each persona, a
    multi-turn dialogue (default 3 turns) runs between a
    \texttt{WikiWriter} (the persona asking questions via DSPy
    \texttt{ChainOfThought}) and a \texttt{TopicExpert} (answering via
    retrieval). Conversations run concurrently.
  \item \textbf{Outline and article generation.} The information table
    from all dialogues feeds hierarchical outline generation (merging a
    ``direct'' and ``informed'' outline), then section-by-section
    article writing with inline citations.
  \item \textbf{Polishing.} A dedicated refinement pass handles
    deduplication and summarisation.
\end{enumerate}

\noindent
STORM uses five separate LLM configurations (conversation, question
asking, outline generation, article writing, polishing), reflecting a
design philosophy where different cognitive tasks benefit from different
model sizes and prompting strategies. Co-STORM~\citep{costorm2024}
extends this with collaborative human-in-the-loop discourse.

The multi-perspective approach is the most architecturally novel
decomposition strategy in this category: by simulating \emph{who would
care about this topic}, it produces more diverse coverage than direct
sub-query generation.

\subsubsection{MiroThinker.}
MiroThinker~\citep{mirothinker2025} (MiroMind AI) achieves the highest
reported BrowseComp score of any system: MiroThinker-H1 reaches 88.2,
surpassing Gemini-3.1-Pro (85.9) and Claude-4.6-Opus (84.0). Limited
architectural details are publicly available; the system is open source
on GitHub.

% =================================================================
\section{Literature Agents}
\label{sec:lit}

Literature agents specialise in answering questions grounded in
academic publications. Unlike deep research systems that search the
open web, these systems operate over curated corpora of scientific
papers, with strong emphasis on citation faithfulness and structured
evidence extraction.

\subsection{PaperQA2}

PaperQA2~\citep{paperqa2024} (FutureHouse) models research as a
reinforcement learning problem. The \texttt{PaperQAEnvironment} class
implements an RL-style \texttt{Environment} with \texttt{reset()} and
\texttt{step(action)} methods, returning observations, rewards, and
termination signals.

The agent selects from five tools at each timestep:

\begin{description}[nosep, style=nextline, labelwidth=3cm, leftmargin=3.5cm]
  \item[\texttt{PaperSearch}] Queries a pre-built search index, adds
    papers to the working \texttt{Docs} collection.
  \item[\texttt{GatherEvidence}] Retrieves relevant chunks from loaded
    papers, scores each chunk 0--10 for relevance via a summary LLM,
    and stores scored \texttt{Context} objects.
  \item[\texttt{GenerateAnswer}] Synthesises an answer from top-ranked
    contexts with inline citations.
  \item[\texttt{Reset}] Clears all gathered evidence (recovery from bad
    search directions).
  \item[\texttt{Complete}] Terminates the session with a certainty
    assessment.
\end{description}

\noindent
Three model roles serve different functions: a main LLM for answer
generation, a summary LLM for evidence scoring, and an embedding model
for semantic search. Paper metadata is enriched via four clients
(Semantic Scholar, CrossRef, OpenAlex, Unpaywall), and a retraction
checker flags withdrawn papers.

On LitQA2, PaperQA2 achieves \textbf{superhuman accuracy}
($\sim$60\% vs.\ $\sim$50\% for PhD-level human experts) with very low
hallucination rates---the strongest evidence that tool-augmented LLM
agents can exceed human performance on academic question answering when
properly grounded in source material.

\subsection{OpenScholar}

OpenScholar~\citep{openscholar2024} (AI2) takes a different approach:
rather than dynamic tool use, it performs dense retrieval over 45M+
full-text paper passages, reranks results, and synthesises answers with
a fine-tuned 8B model specifically trained for citation faithfulness.
On ScholarBench, this 8B model matches or exceeds GPT-4o and PaperQA2
on four scientific QA datasets while using a fraction of the compute.
The key insight is that domain specialisation and retrieval
augmentation can substitute for model scale in well-defined scientific
QA tasks.

\subsection{Elicit}

Elicit (Ought/Elicit Inc.) focuses on \emph{structured extraction}
rather than free-form QA. It searches 200M+ papers via Semantic Scholar,
then uses LLMs to extract structured data columns (study design, sample
size, methods, findings) from abstracts and full texts. This makes it
uniquely suited for systematic reviews and meta-analyses, where the goal
is not a narrative answer but a structured comparison across studies.
It is the most workflow-oriented system in this category, functioning
as an assisted tool rather than an autonomous agent.

\subsection{Consensus}

Consensus indexes 200M+ academic papers and answers questions by
extracting and aggregating claims from abstracts. Its signature feature
is a ``consensus meter'' showing the degree of scientific agreement on
a given claim across studies. A GPT-4-powered synthesis module
(``Consensus Copilot'') provides narrative summaries. The system
demonstrates that useful research automation need not require deep
autonomy---sometimes aggregation across many simple extractions is
more valuable than sophisticated reasoning over a few sources.

% =================================================================
\section{Full-Pipeline Science}
\label{sec:pipeline}

Full-pipeline systems attempt to automate the complete research cycle:
from idea generation through experimentation to paper writing and review.
These are the most ambitious---and most controversial---automated
research systems.

\subsection{The AI Scientist}

The AI Scientist~\citep{aiscientist2024} (Sakana AI) implements a
five-phase sequential pipeline:

\begin{enumerate}[nosep]
  \item \textbf{Idea generation.} The LLM generates research ideas as
    structured JSON (name, title, experiment plan, novelty/feasibility
    ratings), with up to 5 rounds of self-reflection. All previously
    generated ideas are included in the prompt to encourage diversity.
  \item \textbf{Novelty check.} A multi-round interactive loop (up to
    10 iterations) where the LLM searches Semantic Scholar, reviews
    results, and decides whether to search again or declare the idea
    novel/not novel.
  \item \textbf{Experimentation.} Uses Aider (an AI coding assistant) to
    modify \texttt{experiment.py}, run it, and iterate (up to 5 runs,
    4 retries each). The system edits real code and runs real experiments.
  \item \textbf{Writeup.} Section-by-section \LaTeX{} generation with
    per-section prompts, a citation loop (up to 20 rounds of Semantic
    Scholar search and BibTeX integration), and two full refinement
    passes. Compilation errors are auto-corrected via Aider.
  \item \textbf{Review.} Simulated NeurIPS-style peer review with
    ensemble of 5 reviews (temperature 0.1), scoring originality,
    quality, clarity, significance, and soundness.
\end{enumerate}

\noindent
The system produces full \LaTeX{} papers compiled to PDF for $\sim$\$15
per paper (with Claude Sonnet 3.5). Papers score $\sim$4--5/10 on the
automated reviewer---borderline reject at best venues. The primary
limitation is quality: while technically correct, the papers lack
novelty and depth. The system is also restricted to ML/DL domains where
experiments are code-runnable; it cannot perform wet-lab science.

Nevertheless, the AI Scientist demonstrates that end-to-end research
automation is technically feasible, and its open-source release has
enabled significant follow-on work.

\subsection{ShinkaEvolve}

ShinkaEvolve~\citep{shinkaevolve2025} (Sakana AI) shifts from paper
writing to \emph{algorithm discovery} via evolutionary optimisation.
The core loop maintains an archive of evaluated programs, uses LLMs to
generate mutations and crossovers, and evaluates fitness against
concrete metrics. Three innovations drive its sample efficiency:

\begin{description}[nosep, style=nextline, labelwidth=4.5cm, leftmargin=5cm]
  \item[Novelty-based rejection] Before expensive evaluation, program
    text embedding similarity plus an LLM-as-novelty-judge filters
    near-duplicates and trivial variations.
  \item[Multi-armed bandit LLM selection] A bandit dynamically picks the
    best model from an ensemble as evolution progresses, adapting to
    problem structure.
  \item[Exploration/exploitation balancing] An intelligent parent
    sampling strategy decides when to refine known solutions versus try
    novel directions.
\end{description}

\noindent
Results are striking: ShinkaEvolve found a state-of-the-art circle
packing solution ($n=26$) in only 150 samples, where prior systems
(including Google DeepMind's AlphaEvolve) required thousands. It evolved
a 3-stage agentic architecture for AIME math reasoning in 75
generations, and outperformed DeepSeek's Global Load Balancing Loss for
Mixture-of-Experts routing in 30 generations. The system is fully open
source under Apache 2.0, in contrast to AlphaEvolve.

ShinkaEvolve's key insight is that LLM-guided search over program space
is dramatically more sample-efficient than random mutation when combined
with novelty filtering---the LLM's prior knowledge about what
constitutes a ``meaningfully different'' program acts as an implicit
regulariser.

\subsection{Coscientist and ChemCrow}

Coscientist~\citep{coscientist2023} (CMU, published in \emph{Nature})
and ChemCrow~\citep{chemcrow2024} (EPFL) extend automated research into
wet-lab chemistry. Both use tool-augmented LLM agents: ChemCrow
integrates 18 chemistry tools (RDKit, safety databases, reaction
predictors) in a ReAct-style loop; Coscientist additionally controls
robotic lab equipment (Opentrons) to autonomously plan and execute
chemical experiments. Coscientist successfully optimised
palladium-catalysed cross-coupling reactions without human intervention.

These systems demonstrate that the agent paradigm extends beyond
information tasks to physical-world experimentation, though they require
specialised infrastructure and raise safety concerns around autonomous
chemical synthesis.

\subsection{SciAgents}

SciAgents~\citep{sciagents2024} (MIT) takes a unique approach to
hypothesis generation: rather than searching text, it reasons over
\emph{knowledge graphs}. Multiple LLM agents with different roles
(ontologist, scientists, critic) interact over a graph of domain
concepts. The ontology agent structures knowledge; scientist agents
propose hypotheses by finding novel connections in the graph; a critic
agent evaluates plausibility. This graph-based approach is particularly
suited to materials science, where novel material properties often arise
from unexpected combinations of known components.

% =================================================================
\section{Empirical Optimisation}
\label{sec:empirical}

Empirical optimisation systems automate the
\emph{edit--benchmark--iterate} cycle: they modify code, measure the
effect, keep improvements, discard regressions, and repeat. Unlike the
previous categories, these systems generate \emph{new knowledge through
experimentation} rather than synthesising existing knowledge.

\subsection{pi-autoresearch}

pi-autoresearch~\citep{piautoresearch2025} is an extension for the
\texttt{pi} AI coding agent that adds an autonomous experiment loop.
Inspired by Karpathy's autoresearch concept, it targets optimisation
tasks with clear metrics: test speed, bundle size, ML training loss,
build times, Lighthouse scores.

\subsubsection{Architecture.}
The system has three layers:

\begin{description}[nosep, style=nextline, labelwidth=3cm, leftmargin=3.5cm]
  \item[Extension] Registers three tools into the agent's namespace:
    \texttt{init\_experiment} (once), \texttt{run\_experiment}
    (benchmark), and \texttt{log\_experiment} (keep/discard). Manages
    session state and TUI dashboard.
  \item[Skills] Prompt-based instructions (\texttt{SKILL.md} files) that
    teach the agent the full workflow, including how to create benchmark
    scripts, correctness checks, and the loop rules.
  \item[Finalize] A bash script that splits a messy experiment branch
    into clean, independently-mergeable PR branches, verifying that
    their union matches the original.
\end{description}

\subsubsection{Context management.}
pi-autoresearch addresses the context exhaustion problem---the central
challenge for long-running agent sessions---through several mechanisms:

\begin{itemize}[nosep]
  \item \textbf{Append-only JSONL log.} Every experiment result (metric
    values, status, description, confidence) is persisted regardless of
    whether the code change was kept or discarded.
  \item \textbf{ASI (Actionable Side Information).} Structured
    key--value annotations on each experiment that survive code reverts.
    When the agent's context is reset, ASI prevents it from repeating
    failed experiments---the code is gone but the learnings persist.
  \item \textbf{Token-aware auto-resume.} The system tracks token
    consumption per iteration, predicts when the context window will
    overflow (with 1.2$\times$ safety margin), gracefully stops, and
    auto-resumes in a fresh context up to 20 times. This enables
    multi-hour unattended optimisation sessions.
  \item \textbf{Output truncation.} Benchmark output sent to the LLM is
    capped at 10 lines / 4\,KB, with structured \texttt{METRIC
    name=value} lines parsed automatically.
\end{itemize}

\subsubsection{Quality control.}
Confidence scoring uses Median Absolute Deviation (MAD): the ratio
$|\text{best\_delta}| / \text{MAD}$ across all runs distinguishes real
improvements ($\geq$2.0$\times$) from benchmark noise ($<$1.0$\times$).
Optional \emph{backpressure checks} (\texttt{autoresearch.checks.sh})
run correctness validation (tests, types, lint) after each passing
benchmark, preventing optimisation that breaks correctness.

\subsection{Relationship to ShinkaEvolve and AlphaEvolve}

ShinkaEvolve (Section~\ref{sec:pipeline}) operates in the same
edit--evaluate--iterate paradigm but at a different level of
abstraction. Where pi-autoresearch optimises \emph{within a codebase}
(making incremental changes to existing code), ShinkaEvolve evolves
\emph{entire algorithms} (generating novel programs from scratch or via
crossover). AlphaEvolve (Google DeepMind, closed-source) occupies a
similar niche to ShinkaEvolve but requires orders of magnitude more
evaluations.

The key architectural difference is \emph{what persists between
iterations}: pi-autoresearch persists the codebase itself (via git
commits) plus structured learnings (ASI); ShinkaEvolve persists a
population archive with fitness scores; AlphaEvolve persists an
evolutionary island model. All three demonstrate that LLMs are
effective mutation operators for code, but the surrounding
infrastructure---novelty filtering, confidence scoring, context
management---determines practical utility.

% =================================================================
\section{Cross-Cutting Analysis}
\label{sec:cross}

\subsection{Decomposition Strategy}

How a system breaks a research question into sub-tasks is the single
most consequential architectural decision. We observe five distinct
strategies:

\begin{table}[h]
\centering
\small
\begin{tabular}{@{}L{3.5cm}L{3.5cm}L{6cm}@{}}
\toprule
\textbf{Strategy} & \textbf{Systems} & \textbf{Trade-off} \\
\midrule
LLM sub-query generation & GPT-Researcher, OpenAI, Perplexity &
  Simple and fast; queries cluster around obvious angles \\[0.3em]
Multi-perspective simulation & STORM &
  Broader coverage; slower; more creative \\[0.3em]
Interactive novelty search & AI Scientist &
  Good for finding literature gaps; requires domain API \\[0.3em]
Evolutionary mutation & ShinkaEvolve &
  Efficient solution-space exploration; needs fitness function \\[0.3em]
Dynamic agent decision & PaperQA2 &
  Most flexible; hardest to debug and control \\
\bottomrule
\end{tabular}
\caption{Decomposition strategies across systems.}
\label{tab:decomposition}
\end{table}

\subsection{The Speed--Depth Trade-Off}

The central design tension in all automated research systems is how much
compute to spend per query. This manifests differently across categories:

\begin{itemize}[nosep]
  \item \textbf{Deep research:} Perplexity/Grok (1--3\,min, broad
    coverage) vs.\ OpenAI (5--30\,min, deep reasoning).
  \item \textbf{Literature agents:} OpenScholar (single-pass retrieval +
    fine-tuned model) vs.\ PaperQA2 (multi-step tool use).
  \item \textbf{Full-pipeline:} AI Scientist ($\sim$\$15/paper, moderate
    quality) vs.\ human researchers ($\sim$months/paper, high quality).
  \item \textbf{Empirical optimisation:} pi-autoresearch (hours of
    autonomous iteration) vs.\ ShinkaEvolve (150 samples to SOTA).
\end{itemize}

\noindent
No system has resolved this trade-off; the optimal point depends
entirely on the use case. Speed-optimised systems excel at factual
lookups and current-events research; depth-optimised systems excel at
complex analytical questions requiring multi-hop reasoning.

\subsection{Quality Control Mechanisms}

The mechanisms systems use to ensure output quality are the strongest
predictors of practical utility:

\begin{table}[h]
\centering
\small
\begin{tabular}{@{}L{3.8cm}L{3.5cm}L{5.5cm}@{}}
\toprule
\textbf{Mechanism} & \textbf{Systems} & \textbf{Effect} \\
\midrule
Relevance scoring per chunk & PaperQA2 &
  Prevents low-quality evidence from diluting answers \\[0.3em]
Multi-perspective diversity & STORM &
  Prevents tunnel vision and query clustering \\[0.3em]
Novelty rejection sampling & ShinkaEvolve &
  Prevents wasted compute on near-duplicates \\[0.3em]
Backpressure checks & pi-autoresearch &
  Prevents correctness regressions during optimisation \\[0.3em]
Confidence via MAD & pi-autoresearch &
  Distinguishes signal from benchmark noise \\[0.3em]
AI peer review & AI Scientist &
  Catches errors post-hoc (weak calibration) \\[0.3em]
Retraction checking & PaperQA2 &
  Flags withdrawn papers before citation \\
\bottomrule
\end{tabular}
\caption{Quality control mechanisms and their effects.}
\label{tab:quality}
\end{table}

\subsection{Context Management}

Long-running agent sessions inevitably exhaust context windows. Systems
address this in three ways:

\begin{enumerate}[nosep]
  \item \textbf{Large windows.} Commercial systems use 128K--1M token
    contexts (Gemini, Claude). Effective but expensive, and still
    insufficient for multi-hour sessions.
  \item \textbf{Compression.} GPT-Researcher caps deep research at
    25K words; PaperQA2 uses relevance scoring to prioritise what enters
    context.
  \item \textbf{External memory.} pi-autoresearch uses append-only logs,
    ASI annotations, and session documents to persist learnings across
    context resets. This is the most sophisticated approach and the only
    one that enables truly multi-hour autonomous sessions.
\end{enumerate}

\subsection{Multi-Model Architectures}

Most systems use multiple LLM configurations for different cognitive
tasks. STORM uses five (conversation, question asking, outline, article
writing, polishing). GPT-Researcher uses two (strategic for planning,
smart for writing). PaperQA2 uses three (main, summary, embedding).
ShinkaEvolve dynamically selects from an ensemble via multi-armed
bandit.

The pattern suggests that \emph{no single model is optimal for all
sub-tasks} in a research workflow: lightweight models suffice for
evidence scoring and query generation, while capable models are needed
for synthesis and reasoning.

\subsection{Open Source vs.\ Commercial}

Open-source systems have reached or exceeded commercial systems on
key benchmarks:

\begin{itemize}[nosep]
  \item \textbf{BrowseComp:} MiroThinker-H1 (88.2) $>$ all commercial
    systems.
  \item \textbf{Algorithm discovery:} ShinkaEvolve (150 samples) beats
    AlphaEvolve (thousands of samples).
  \item \textbf{Scientific QA:} OpenScholar 8B matches GPT-4o on
    ScholarBench.
  \item \textbf{Literature QA:} PaperQA2 exceeds human PhD-level
    performance on LitQA2.
\end{itemize}

\noindent
The commercial advantage lies in integration (search quality, UI,
ecosystem) rather than raw capability. For researchers building custom
systems, open-source components provide building blocks that match or
exceed proprietary alternatives.

% =================================================================
\section{Conclusion}
\label{sec:conclusion}

The automated research landscape has matured rapidly from single-turn
RAG systems to multi-agent architectures capable of extended autonomous
research sessions. Our four-category taxonomy---deep research, literature
agents, full-pipeline science, and empirical optimisation---captures
fundamentally different automation targets that share common
architectural patterns but diverge in decomposition strategy, quality
control, and context management.

Three findings stand out:

\begin{enumerate}[nosep]
  \item \textbf{Decomposition matters most.} STORM's multi-perspective
    simulation produces more diverse research than direct sub-query
    generation; ShinkaEvolve's novelty-filtered evolution is dramatically
    more sample-efficient than random mutation. How a system breaks
    problems apart determines coverage, efficiency, and output quality.
  \item \textbf{Quality control is the differentiator.} The gap between
    a useful and useless research agent is not the underlying LLM but
    the surrounding quality infrastructure: relevance scoring (PaperQA2),
    confidence estimation (pi-autoresearch), novelty filtering
    (ShinkaEvolve), and perspective diversity (STORM).
  \item \textbf{The two paradigms will converge.} Information synthesis
    (finding existing knowledge) and empirical discovery (generating new
    knowledge through experimentation) are currently separate. The most
    impactful future systems will combine both---using literature agents
    to generate hypotheses, then empirical optimisation to test them.
    FutureHouse Robin and the AI Scientist are early attempts, but
    neither does both halves well yet.
\end{enumerate}

\noindent
For practitioners building automated research systems, the evidence
favours composing specialised open-source components (PaperQA2 for
literature, STORM-style decomposition for breadth, pi-autoresearch-style
context management for long sessions) over relying on monolithic
commercial offerings. The tools are available; the remaining challenge is
integration.

% =================================================================
\bibliographystyle{plainnat}

\begin{thebibliography}{99}

\bibitem[{Boiko et~al.(2023)}]{coscientist2023}
Boiko, D.~A., MacKnight, R., and Gomes, G.
\newblock Autonomous chemical research with large language models.
\newblock \emph{Nature}, 624:570--578, 2023.

\bibitem[{Bran et~al.(2024)}]{chemcrow2024}
Bran, A.~M., Cox, S., Schilter, O., Baldassari, C., White, A.~D., and
Schwaller, P.
\newblock ChemCrow: Augmenting large-language models with chemistry tools.
\newblock \emph{Nature Machine Intelligence}, 6:525--535, 2024.

\bibitem[{Elovic(2024)}]{gptresearcher2024}
Elovic, A.
\newblock GPT-Researcher: Autonomous agent for comprehensive online research.
\newblock \url{https://github.com/assafelovic/gpt-researcher}, 2024.

\bibitem[{Lu et~al.(2024)Lu, Lu, Lange, Foerster, Clune, and
  Ha}]{aiscientist2024}
Lu, C., Lu, Y., Lange, R.~T., Foerster, J., Clune, J., and Ha, D.
\newblock The {AI Scientist}: Towards fully automated open-ended scientific
  discovery.
\newblock \emph{arXiv preprint arXiv:2408.06292}, 2024.

\bibitem[{Lu et~al.(2024)Lu, Choi}]{storm2024}
Lu, Y., et~al.
\newblock {STORM}: Synthesis of topic outlines through retrieval and
  multi-perspective question asking.
\newblock \emph{arXiv preprint arXiv:2402.14207}, 2024.

\bibitem[{Jiang et~al.(2024)}]{costorm2024}
Jiang, Y., et~al.
\newblock Into the unknown unknowns: Engaged human learning through
  participation in language model agent conversations.
\newblock \emph{arXiv preprint arXiv:2408.15232}, 2024.

\bibitem[{Lala et~al.(2024)}]{paperqa2024}
Lala, J., O'Donoghue, O., Mondal, A., Wallace, B., and White, A.~D.
\newblock {PaperQA}: Retrieval-augmented generative agent for scientific
  research.
\newblock \emph{arXiv preprint arXiv:2312.07559}, 2024.

\bibitem[{Asai et~al.(2024)}]{openscholar2024}
Asai, A., et~al.
\newblock {OpenScholar}: Synthesizing scientific literature with
  retrieval-augmented {LM}s.
\newblock \emph{arXiv preprint arXiv:2411.14199}, 2024.

\bibitem[{Buehler(2024)}]{sciagents2024}
Buehler, M.~J., et~al.
\newblock {SciAgents}: Automating scientific discovery through multi-agent
  intelligent graph reasoning.
\newblock \emph{arXiv preprint arXiv:2409.05556}, 2024.

\bibitem[{Sakana AI(2025)}]{shinkaevolve2025}
Sakana AI.
\newblock {ShinkaEvolve}: Evolutionary code optimization with {LLMs}.
\newblock \emph{arXiv preprint arXiv:2509.19349}, 2025.

\bibitem[{Zechner and Beorlegui(2025)}]{piautoresearch2025}
Zechner, M. and Beorlegui, D.
\newblock pi-autoresearch: Autonomous experiment loops for {AI} coding agents.
\newblock \url{https://github.com/davebcn87/pi-autoresearch}, 2025.

\bibitem[{OpenAI(2025)}]{browsecomp2025}
OpenAI.
\newblock {BrowseComp}: A simple challenge for browsing agents.
\newblock \url{https://openai.com/index/browsecomp/}, 2025.

\bibitem[{MiroMind AI(2025)}]{mirothinker2025}
MiroMind AI.
\newblock {MiroThinker}: Open research agent.
\newblock \url{https://github.com/MiroMindAI/MiroThinker}, 2025.

\bibitem[{Perplexity AI(2026)}]{draco2026}
Perplexity AI.
\newblock {DRACO}: Evaluating deep research performance in the wild.
\newblock \emph{arXiv preprint arXiv:2602.11685}, 2026.

\end{thebibliography}

\end{document}
